% This document is compiled with LuaLaTeX to support Korean in form fields.
% It uses the 'kotex' package for Korean language support.
\documentclass[10pt, a4paper, landscape]{article}

% 1. PACKAGES
% \usepackage[T1]{fontenc}  % REMOVED: Not needed for LuaLaTeX
% \usepackage{lmodern}      % REMOVED: Not needed for LuaLaTeX
\usepackage[landscape, margin=1cm]{geometry} % Page layout (landscape, small margins)
\usepackage{array}          % For \extrarowheight
\usepackage[unicode]{hyperref} % MOVED: Load hyperref *before* kotex for font setup
\usepackage{kotex}          % For Korean language support
\usepackage{forloop}        % For looping pages




% 2. FORM FIELD DEFINITION
% Define a custom command \tf for text fields
% Usage: \tf{<base_name>}{<width>}{<page_id>}
% This creates a standard, bordered text field with a unique name.
\newcommand{\tf}[3]{%
  \TextField[name=#1_#3, width=#2, height=10pt, charsize=8pt, bordercolor={0.5 0.5 0.5}]{}%
}

% We define a command \tfrow for a single data row to avoid repetition
% \tfrow{<row_id>}{<page_id>}
\newcommand{\tfrow}[2]{
  \tf{date_#1}{2.0cm}{#2} & \tf{item_#1}{3.0cm}{#2} & \tf{spec_#1}{2.0cm}{#2} & \tf{qty_#1}{2.0cm}{#2} & \tf{unit_#1}{2.0cm}{#2} & \tf{supply_#1}{2.6cm}{#2} & \tf{tax_#1}{2.6cm}{#2}
}

% 3. DOCUMENT BODY
\begin{document}

% Create a counter for our page loop
\newcounter{loopcounter}

% The \begin{Form} environment contains all form fields
\begin{Form}[NeedAppearances=true]

% Set table properties
\setlength{\extrarowheight}{12pt} % Add padding to rows
\setlength{\tabcolsep}{4pt}       % Reduce space between columns

% --- PAGE LOOP START ---
% We loop from 1 to 3 (n=3 pages). 
% To change the number of pages, change the '4' (which means 'less than 4').
% For 5 pages, use: \value{loopcounter} < 6
\forloop{loopcounter}{1}{\value{loopcounter} < 11}{%

% --- Add visible page number ---
\noindent % Prevent indentation
\parbox{\linewidth}{\raggedleft Page \theloopcounter}
\vspace{2pt} % Add a little space before the table

% We use \noindent to prevent indentation before the table
\noindent
% We use a 'tabular' environment with manually defined column widths
% to precisely match the layout of the image.
% Total 7 columns:
% C1: 2.2cm | C2: 3.2cm | C3: 2.2cm | C4: 2.2cm | C5: 2.2cm | C6: 2.8cm | C7: 2.8cm
\begin{tabular}{|p{2.2cm}|p{3.2cm}|p{2.2cm}|p{2.2cm}|p{2.2cm}|p{2.8cm}|p{2.8cm}|}
\hline

% --- ROW 1: Date and Registration No. ---
% Spans C1-C3
\multicolumn{3}{|p{7.6cm}|}{ % 2.2+3.2+2.2 + 4*tabcolsep
  \tf{year}{1.2cm}{\theloopcounter} 년 \tf{month}{1.2cm}{\theloopcounter} 월 \tf{day}{1.2cm}{\theloopcounter} 일
} &
% Spans C4-C5
\multicolumn{2}{c|}{등록번호} &
% Spans C6-C7
\multicolumn{2}{p{5.6cm}|}{ \tf{reg_no_1}{5.4cm}{\theloopcounter} } \\
\hline

% --- ROW 2: Recipient, Company, Name ---
\multicolumn{3}{|c|}{ \tf{attn_name}{3.6cm}{\theloopcounter} 귀하 } & % C1
\multicolumn{1}{c|}{상호} & % C2
\multicolumn{1}{p{2.2cm}|}{ \tf{sender_co}{2.0cm}{\theloopcounter} } & % C3
\multicolumn{1}{c|}{성명} & % C4
% Spans C5-C7
\multicolumn{1}{p{2.6cm}|}{ \tf{sender_name}{2.6cm}{\theloopcounter} } \\
\hline

% --- ROW 3: Calculation note, Address ---
% Spans C1-C3
\multicolumn{3}{|c|}{아래와 같이 계산합니다} &
% Spans C4-C5
\multicolumn{2}{c|}{사업장 주소} &
% Spans C6-C7
\multicolumn{2}{p{5.6cm}|}{ \tf{sender_addr}{5.4cm}{\theloopcounter} } \\
\hline

% --- ROW 4: Biz Type, Biz Item ---
\multicolumn{3}{|c|}{아래와 같이 계산합니다} & % C1-C3 (Empty)
\multicolumn{1}{c|}{업태} & % C4
\multicolumn{1}{p{2.2cm}|}{ \tf{sender_biz_type}{2.0cm}{\theloopcounter} } & % C5
\multicolumn{1}{c|}{종목} & % C6
\multicolumn{1}{p{2.8cm}|}{ \tf{sender_biz_item}{2.6cm}{\theloopcounter} } \\ % C7
\hline

% --- ROW 5: Total Amount, Market Price header ---
\multicolumn{1}{|c|}{합계금액} & % C1
% Spans C2-C3
\multicolumn{2}{p{5.4cm}|}{ \tf{total_amount_text}{5.2cm}{\theloopcounter} } &
% Spans C4-C7
\multicolumn{1}{c|}{시세} & 
\multicolumn{3}{p{5.4cm}|}{ \tf{total_amount}{5.2cm}{\theloopcounter} }  \\
\hline

% --- ROW 6: Main Table Headers ---
\multicolumn{1}{|c|}{월일} &     % C1
\multicolumn{1}{c|}{품목} &     % C2
\multicolumn{1}{c|}{규격} &     % C3
\multicolumn{1}{c|}{수량} &     % C4
\multicolumn{1}{c|}{단가} &     % C5
\multicolumn{1}{c|}{공급가액} & % C6
\multicolumn{1}{c|}{세액} \\      % C7
\hline

% --- Data Rows (x12) ---
% Generate the 12 data rows
\tfrow{1}{\theloopcounter} \\ \hline
\tfrow{2}{\theloopcounter} \\ \hline
\tfrow{3}{\theloopcounter} \\ \hline
\tfrow{4}{\theloopcounter} \\ \hline
\tfrow{5}{\theloopcounter} \\ \hline
\tfrow{6}{\theloopcounter} \\ \hline
\tfrow{7}{\theloopcounter} \\ \hline
\tfrow{8}{\theloopcounter} \\ \hline
\tfrow{9}{\theloopcounter} \\ \hline
\tfrow{10}{\theloopcounter} \\ \hline
\tfrow{11}{\theloopcounter} \\ \hline
\tfrow{12}{\theloopcounter} \\ \hline

% --- Subtotal Row ---
% Spans C1-C5
\multicolumn{3}{|c|}{합계} &
\multicolumn{3}{|c|}{\tf{total_supply}{6cm}{\theloopcounter}} &
\tf{total_tax}{2.6cm}{\theloopcounter} \\   % C7
\hline

% --- Footer Row ---
\multicolumn{1}{|c|}{입금} & % C1
\multicolumn{1}{p{3.2cm}|}{ \tf{f_deposit}{3.0cm}{\theloopcounter} } & % C2
\multicolumn{1}{c|}{잔금} & % C3
\multicolumn{1}{p{2.2cm}|}{ \tf{f_balance}{2.0cm}{\theloopcounter} } & % C4
\multicolumn{1}{c|}{연수자} & % C5
% Spans C6-C7
\multicolumn{2}{p{5.6cm}|}{ \tf{f_collector}{5.4cm}{\theloopcounter} } \\
\hline

\end{tabular} % End of the table

\clearpage % Force a new page for the next loop iteration

}% --- PAGE LOOP END ---

\end{Form} % End of the form
\end{document}

% End of document

% LuaLaTeX compilation command:
% cd etc && 
% lualatex -shell-escape -synctex=1 -interaction=nonstopmode form.tex
% Note: Ensure that your LaTeX editor is set to compile with LuaLaTeX to properly render Korean characters in form fields.